%%%%%%%%%%%%%%%%%%%%%%%%%%%%%%%%%
%
%       EXERCÍCIO
%
%%%%%%%%%%%%%%%%%%%%%%%%%%%%%%%%%

\ifdefstring{\atividade}{prova}{%
    \renewcommand{\valorquestao}{\ValorQ\ pontos}
}%

\begin{exercicioBanco}[\valorquestao]
Um caminho para chegar a uma festa pode ser dividido em três etapas. Sem enganos, o trajeto é feito em 1 hora e 30 minutos. Se enganos acontecem na primeira etapa, acrescente 15 minutos ao tempo do trajeto. Para enganos na segunda etapa, o acréscimo é de 20 minutos e, para a terceira, 25 minutos. Admita que a probabilidade de engano é 0,15; 0,20 e 0,25 para a primeira, segunda e terceira etapas, respectivamente. É provável haver atraso na chegada à festa? Determine a probabilidade de haver atraso e o atraso não passar de 40 minutos.
\end{exercicioBanco}

